\section*{Hoofdstuk \ref{ch:b1:alcohol-activation} --- Alcoholen: de OH-groep activeren}

\begin{solution}{exo:b1:alcohol-activation:1}
\ce{2CH3CH2OH + 2Na -> 2CH3CH2O- + 2Na+ + H2};
\ce{CH3CH2OH + NaH -> CH3CH2O- + Na+ + H2};
\ce{CH3CH2OH + OH- <=> CH3CH2O- + H2O}, $K = 10^{14.0 - 15.9} = 10^{-1.9}$:
gedeeltelijk.
\end{solution}

\begin{solution}{exo:b1:alcohol-activation:2}
1-Methoxypropaan; methoxyethaan; ethoxybenzeen (het fenoxide, een
zwakkere base dan een \omterm{def:b1:alcohol-activation:alkoxide}{alkoxide}, is nog steeds een goed \omterm{def:b1:electronic-effects:nucleophile}{nucleofiel}).
\end{solution}

\begin{solution}{exo:b1:alcohol-activation:3}
Tosylering in pyridine geeft \ce{CH3CH2CH2OTs}; daarna
\[ \ce{CH3CH2CH2OTs + I- -> CH3CH2CH2I + TsO-}. \] Rechtstreeks zou jodide
\ce{OH-} moeten verdrijven: geen reactie. Het \omterm{def:b1:alcohol-activation:sulfonate}{tosylaat} heeft een
uitstekende \omterm{def:b1:electronic-effects:nucleophile}{vertrekkende groep}, en de omstandigheden zijn zuur noch
sterk basisch.
\end{solution}

\begin{solution}{exo:b1:alcohol-activation:4}
Butaan-2-ol: but-1-een, \cip{Z}- en \cip{E}-but-2-een; hoofdproduct
\cip{E}-but-2-een. 2-Methylpentaan-2-ol: 2-methylpent-2-een
(hoofdproduct, trigesubstitueerd) en 2-methylpent-1-een.
\end{solution}

\begin{solution}{exo:b1:alcohol-activation:5}
\ce{(CH3)2CHCH2-O-CH2CH3}: natrium-2-methylpropaan-1-olaat met
broomethaan, of natriumethoxide met 1-broom-2-methylpropaan. Beide
halogeniden zijn primair, maar 1-broom-2-methylpropaan is vertakt naast
het reagerende koolstofatoom, wat SN2 vertraagt en de \omterm{def:b1:predominant-reaction:strong-acid}{sterke base} laat
elimineren (2-methylpropeen). De eerste keuze is beter.
\end{solution}

\begin{solution}{exo:b1:alcohol-activation:6}
Het \omterm{def:b1:alcohol-activation:sulfonate}{tosylaat} behoudt de \cip{R}-configuratie; SN2 met ethoxide keert haar
om: \cip{S}-2-ethoxyoctaan (OEt krijgt de eerste rang, zoals OTs). Met
zwavelzuur en warmte: eliminatie tot octenen (vooral
\cip{E}-oct-2-een), via een secundair \omterm{def:b1:electronic-effects:intermediates}{carbokation}, met verlies van de
stereochemie.
\end{solution}

\begin{solution}{exo:b1:alcohol-activation:7}
Protonering van OH; verlies van water tot het tertiaire \omterm{def:b1:electronic-effects:intermediates}{carbokation}
(traag); invangen door \ce{Cl-}: 2-chloor-2-methylpropaan. Snel omdat het
tertiaire \omterm{def:b1:electronic-effects:intermediates}{carbokation} stabiel is en het chloride, onoplosbaar in het
zuur, zich afscheidt.
\end{solution}

\begin{solution}{exo:b1:alcohol-activation:8}
\ce{CH3CH2OH + H+ -> CH3CH2OH2+}; een tweede ethanolmolecuul valt het
koolstofatoom van achteren aan terwijl water vertrekt (SN2):
\ce{(CH3CH2)2OH+}, dat een proton verliest: ethoxyethaan. Concurrerend:
een base (ethanol, \ce{HSO4-}) neemt een $\beta$-proton van
\ce{CH3CH2OH2+} terwijl water vertrekt (E2): etheen.
\end{solution}

\begin{solution}{exo:b1:alcohol-activation:9}
2-Methylpropaan-2-ol (tertiair), dat het onoplosbare
2-chloor-2-methylpropaan geeft.
\end{solution}

\begin{solution}{exo:b1:alcohol-activation:10}
Protonering en verlies van water geven een secundair \omterm{def:b1:electronic-effects:intermediates}{carbokation} naast
een quaternair koolstofatoom. Een methylgroep van dat koolstofatoom
verschuift met haar \omterm{def:b1:lewis-resonance-vsepr:lewis}{bindend paar} naar het kationische koolstofatoom: de
positieve lading zit nu op een \omterm{def:b1:electronic-effects:carbon-class}{tertiair koolstofatoom}, stabieler. Verlies
van een proton van een naburig koolstofatoom geeft het
tetragesubstitueerde 2,3-dimethylbut-2-een, het stabielste alkeen.
\end{solution}

\begin{solution}{exo:b1:alcohol-activation:11}
Uit \cip{R}-butaan-2-ol: \omterm{def:b1:alcohol-activation:sulfonate}{tosylaat} (\cip{R}), daarna natriummethoxide
(SN2, inversie): \cip{S}-2-methoxybutaan. Uit \cip{S}-butaan-2-ol: maak
het \omterm{def:b1:alcohol-activation:alkoxide}{alkoxide} (NaH) en voeg joodmethaan toe: de stereogene C--O-binding
blijft onaangeroerd, opnieuw \cip{S}-2-methoxybutaan.
\end{solution}

\begin{solution}{exo:b1:alcohol-activation:12}
$Q = [\text{ether}][\ce{H2O}]/[\ce{EtOH}]^2$. Water verwijderen naarmate
het ontstaat, houdt $[\ce{H2O}]$ en $Q$ klein: $Q < K$ en de reactie blijft
ether vormen tot de alcohol verbruikt is.
\end{solution}

\begin{solution}{pb:b1:alcohol-activation:1}
\textbf{1.} Natrium-tert-butoxide met een methylhalogenide;
natriummethoxide met een tert-butylhalogenide.
\textbf{2.} Het tweede: methoxide, een \omterm{def:b1:predominant-reaction:strong-acid}{sterke base}, elimineert uit het
tertiaire halogenide, \ce{(CH3)3CBr + CH3O- -> (CH3)2C=CH2 + CH3OH + Br-}.
\textbf{3.} \ce{(CH3)3CO- + CH3I -> (CH3)3COCH3 + I-}: de zuurstof van het
\omterm{def:b1:alcohol-activation:alkoxide}{alkoxide} valt het methylkoolstofatoom van achteren aan terwijl jodide
vertrekt (SN2).
\textbf{4.} Met natrium (of natriumhydride): \ce{2(CH3)3COH + 2Na ->
2(CH3)3CO- + 2Na+ + H2}. Hydroxide is een te \omterm{def:b1:predominant-reaction:strong-acid}{zwakke base}: $K = 10^{14.0 -
19.2} = 10^{-5.2}$.
\textbf{5.} Water zou het \omterm{def:b1:alcohol-activation:alkoxide}{alkoxide} protoneren en het halogenide aanvallen.
\textbf{6.} Nee: geen \omterm{def:b1:stereochemistry-in-depth:stereocentre}{stereogeen centrum}.
\textbf{7.} \cip{R}-butaan-2-ol + TsCl (pyridine) $\to$
\cip{R}-butaan-2-yltosylaat: retentie.
\textbf{8.} \cip{S}-2-methoxybutaan, door SN2-inversie.
\textbf{9.} \cip{S}-butaan-2-ol.
\textbf{10.} \cip{R}-2-methoxybutaan: de C--O-binding van het stereogene
koolstofatoom blijft behouden; alleen O--C(methyl) wordt gevormd.
\textbf{11.} E2 (methoxide is een \omterm{def:b1:predominant-reaction:strong-acid}{sterke base} en het koolstofatoom is
secundair), wat butenen geeft; beperkt door de kleine base, een
uitstekende \omterm{def:b1:electronic-effects:nucleophile}{vertrekkende groep} en een milde temperatuur.
\textbf{12.} De geprotoneerde alcohol ioniseert tot een vlak secundair
\omterm{def:b1:electronic-effects:intermediates}{carbokation}, dat aan beide zijden door methanol ingevangen wordt:
racemische ether, samen met butenen.
\textbf{13.} Protonering van OH, verlies van water (traag) tot het
tertiaire \omterm{def:b1:electronic-effects:intermediates}{carbokation}, verlies van een $\beta$-proton.
\textbf{14.} 2-Methylbut-2-een (hoofdproduct, trigesubstitueerd) en
2-methylbut-1-een.
\textbf{15.} Een tertiair \omterm{def:b1:electronic-effects:intermediates}{carbokation} is gehinderd en verliest veel
gemakkelijker een proton dan dat een ander alcoholmolecuul het kan
aanvallen.
\textbf{16.} Om een product te verwijderen: zolang $Q$ klein blijft, gaat
de reactie door, en het alkeen blijft niet in contact met het zuur.
\textbf{17.} Een hoge eerste barrière (verlies van water,
snelheidsbepalend) naar het \omterm{def:b1:electronic-effects:intermediates}{carbokation}, daarna een kleine barrière naar
het alkeen.
\textbf{18.} Nee: een primair \omterm{def:b1:electronic-effects:intermediates}{carbokation} is te instabiel; primaire
alcoholen elimineren via E2 op de geprotoneerde alcohol, bij hogere
temperatuur.
\textbf{19.} $10.0/74.0 = \qty{0.135}{mol}$.
\textbf{20.} $1.10 \times 0.135 \times 141.9 = \qty{21.1}{g}$.
\textbf{21.} $0.135 \times 88.0 = \qty{11.9}{g}$.
\textbf{22.} $0.75 \times 11.9 =$ \textbf{\qty{8.9}{g} MTBE}.
\end{solution}
